\chapter{Componentes}
\label{ch:componentes}

En este capítulo se enumeran y describen los componentes utilizados en el prototipo, junto con la justificación técnica de su selección, los parámetros relevantes para el modelo y las consideraciones de integración con \gls{ROS2}, montaje y seguridad. El robot corresponde a una plataforma diferencial con dos ruedas motrices y una rueda loca (\emph{caster}), un escáner láser \gls{LIDAR} RPLidar~A1, una \gls{RaspberryPi}~4 como computador de a bordo y un \gls{Arduino}~Nano para el control de bajo nivel.

%------------------------------------------------------------------
\section{Arquitectura física del robot}

\subsection{Plataforma diferencial y chasis}

El prototipo adopta una arquitectura de tracción diferencial con dos ruedas motrices independientes y una rueda loca de apoyo. La pertinencia de esta configuración para la navegación de interiores de bajo costo —maniobrabilidad, calidad de odometría y madurez del soporte en \gls{ROS2}— se argumentó en el Capítulo~\ref{ch:estado_arte}, y su modelo cinemático se desarrolló en el Capítulo~\ref{marco_teorico}. Esta sección documenta la realización física de dicha arquitectura y los parámetros geométricos que alimentan el modelo de odometría: el radio de rueda $r$ y la distancia entre ruedas $L$.

\utemFigure{2.Texto/Imag/diff_bec.png}{0.40\textwidth}
{Arquitectura móvil de tracción diferencial: dos ruedas motrices y una rueda loca (caster).}
{BECChassisLong}{fig:chasis-2wd-caster}

La rueda loca aporta estabilidad sin influir en la dirección, de modo que el control se ejerce íntegramente mediante la diferencia de velocidad entre las dos ruedas motrices. Las alternativas Ackermann y omnidireccional, junto con la comparación cualitativa que respalda la elección de la arquitectura diferencial, se documentan en el Anexo~\ref{anx:arquitecturas}.

\subsection{Ruedas motrices}

Se seleccionaron ruedas estándar de goma para la tracción, por su robustez, economía y desempeño consistente en interiores. Su dibujo antideslizante incrementa la adherencia en superficies comunes como baldosa, vinilo o madera.

\utemFigure{2.Texto/Imag/wheel_set.png}{0.50\textwidth}
{Conjunto de ruedas.}
{ThinkRobotics65mmWheel}{fig:wheel_set}

\begin{table}[H]
\centering
\begin{threeparttable}
  \caption{Especificaciones geométricas de las ruedas motrices.}
  \label{tab:ruedas-especificaciones}
  \small
  \begin{tabular}{l S[table-format=3.2] l l}
    \toprule
    \textbf{Parámetro} & {\textbf{Valor (mm)}} & \textbf{Valor (m)} & \textbf{Observación} \\ 
    \midrule
    Diámetro ($d$)      & 67.00  & 0.06700 & Dato del fabricante \\
    Ancho               & 26.00  & 0.02600 & Dato del fabricante \\ \addlinespace
    Radio ($r$)         & 33.50  & 0.03350 & Derivado ($r = d/2$) \\
    Circunferencia ($C$) & 210.49 & 0.21049 & Derivado ($C = \pi d$) \\ 
    \bottomrule
  \end{tabular}
  \begin{tablenotes}
    \footnotesize
    \item[] El radio $r$ y la circunferencia $C$ alimentan el cálculo de odometría del Capítulo~\ref{marco_teorico}.
  \end{tablenotes}
\end{threeparttable}
\end{table}

\subsection{Rueda loca (\emph{caster})}

Se emplea un rodillo esférico tipo \emph{ball caster} como rueda de apoyo. Su esfera metálica gira libremente en todas direcciones, aportando estabilidad sin influir en la dirección. Sus dimensiones reducidas (2.8~×~4.5~cm) mantienen el punto de apoyo trasero cercano al suelo, favoreciendo un centro de gravedad bajo, y soporta cargas muy por encima del peso estimado del robot.

\utemFigure{2.Texto/Imag/ball_caster.png}{0.35\textwidth}{Rueda loca.}
{FIXOhu_OnmagarasztoGorgo}{fig:ball_caster}

%------------------------------------------------------------------
\section{Actuación: motores DC con caja reductora}

Se utiliza el motorreductor JGA25-370 (12~V, 130~RPM), seleccionado por su tensión nominal compatible (12~V), su par suficiente para superar el rozamiento estático y pequeñas pendientes, su corriente dentro del margen del driver L298N, su bajo \emph{backlash} y su caja reductora integrada (relación aproximada 1:45), que convierte alta velocidad y bajo par en una salida de menor velocidad y mayor par, adecuada para tracción. La relación de reducción $G\approx45$ interviene directamente en el cálculo de la resolución del encoder (Sección~\ref{sec:encoder}).

\utemFigure{2.Texto/Imag/motor.jpg}{0.5\linewidth}{Motor modelo JGA25-370, 12~V, 130~RPM.}{nfp_gm25_370_en}{fig:motor}

\paragraph{Especificaciones técnicas relevantes}
\begin{itemize}
  \item \textbf{Tensión nominal}: 12~V DC.
  \item \textbf{Velocidad sin carga}: 130~RPM.
  \item \textbf{Corriente sin carga}: 0.05--0.15~A.
  \item \textbf{Ratio de reducción}: aproximadamente 1:45.
  \item \textbf{Torque nominal}: $\sim$1~kg$\cdot$cm; torque de bloqueo $\sim$3.6~kg$\cdot$cm a 2~A.
  \item \textbf{Peso}: $\sim$110~g.
\end{itemize}

%------------------------------------------------------------------
\section{Encoder de cuadratura}
\label{sec:encoder}

Se selecciona un encoder incremental en cuadratura de efecto Hall integrado a la familia JGA25. Sus ventajas son la compatibilidad mecánica directa con el motor (montaje en la tapa trasera, sin soportes externos), la robustez frente a polvo y vibraciones (superior a los encoders ópticos de bajo costo), la señal digital A/B en cuadratura a 90° —que permite determinar dirección y velocidad— y la compatibilidad eléctrica a 5~V, con conexión directa al \gls{Arduino}.

\utemFigure{2.Texto/Imag/encoder.jpg}{0.7\linewidth}{Conexiones del encoder incremental en cuadratura de efecto Hall.}{gm25_370_encoder}{fig:encoder}

\subsection{Resolución efectiva}

La resolución en el eje de salida (rueda) se calcula como:

\begin{equation}
\mathrm{CPR}_{\text{salida}} = N_{\mathrm{Hall}} \cdot s \cdot G
\end{equation}

donde $N_{\mathrm{Hall}}=11$ pulsos por vuelta del motor, $s=4$ (decodificación en cuadratura completa) y $G\approx45$ (relación de reducción). Así:

\begin{equation}
\mathrm{CPR}_{\text{salida}} = 11 \times 4 \times 45 = 1980 \text{ cuentas/vuelta}
\end{equation}

Con rueda de diámetro $d=67\,\text{mm}$ ($C\approx0.2105\,\text{m}$), la resolución lineal es $\approx 1980 / 0.2105 \approx 9.41 \times 10^3$ cuentas/m. Este valor ($N_e = 1980$) corresponde a la resolución del encoder empleada en las ecuaciones de odometría del Capítulo~\ref{marco_teorico} y en la calibración del firmware.

%------------------------------------------------------------------
\section{Driver de motores L298N}

Se adopta el módulo L298N como puente H dual por su disponibilidad, bajo costo y compatibilidad lógica a 5~V con el \gls{Arduino} \citep{STMicro_L298N_Datasheet}. Sus principales especificaciones son:

\begin{itemize}
  \item Topología: puente H dual para motores DC bidireccionales.
  \item Alimentación de motores: el módulo \emph{breakout} típico opera entre 7--35~V debido al regulador 78M05 integrado.
  \item Alimentación lógica: 4.5--7~V (típico 5~V, compatible TTL).
  \item Corriente por canal: 2.0~A continua; 3.0~A pico no repetitivo ($t \leq 100$~ms).
  \item Control de velocidad mediante \gls{PWM} en pines ENA y ENB.
\end{itemize}

La zona muerta del L298N constituye una de las no linealidades que el refinamiento experimental del controlador PID compensa, tal como se anticipó en el Capítulo~\ref{marco_teorico}.

%------------------------------------------------------------------
\section{Microcontrolador de bajo nivel: Arduino Nano V3}

Se emplea un \gls{Arduino}~Nano~V3 (ATmega328P) para ejecutar las tareas sensibles al tiempo: generación de PWM, lectura de encoders por interrupciones y control PID de velocidad. De esta forma se aísla el lazo de control del \emph{jitter} del sistema operativo de la Raspberry Pi: el PID corre a 30~Hz en el microcontrolador mientras la Raspberry Pi se dedica a \gls{ROS2}, SLAM y navegación.

\paragraph{Especificaciones técnicas} \citep{ArduinoNanoStore,ArduinoDocsNano,Atmega328PDatasheet}
\begin{itemize}
  \item Microcontrolador: ATmega328P a 16~MHz, lógica de 5~V.
  \item Memoria: 32~KB Flash, 2~KB SRAM.
  \item E/S: 14 digitales (6 PWM), 8 analógicas (10 bits).
  \item Interrupciones externas en D2 y D3, y por cambio de pin (PCINT).
\end{itemize}

\subsection{Justificación e integración}

El \gls{Arduino}~Nano se eligió por su compatibilidad eléctrica a 5~V con el L298N, su tamaño compacto y su ecosistema maduro. Su microcontrolador ATmega328P de 8 bits (2~KB de SRAM) queda por debajo de los requisitos mínimos de \gls{microRos}~\cite{microros} —el \textit{framework} de \gls{ROS2} para microcontroladores de 32 bits—, por lo que la integración con el grafo de \gls{ROS2} se resolvió mediante un puente serial alojado en la Raspberry Pi (Sección~\ref{sec:integracion_ros2}). Plataformas como ESP32, Teensy~4.x o Raspberry Pi Pico (RP2040) sí cumplen esos requisitos y constituyen candidatas naturales para una migración futura que conserve la misma arquitectura mecánica y de control.

%------------------------------------------------------------------
\section{Raspberry Pi 4 (8 GB)}

La Raspberry Pi~4 (8~GB) \parencite{RaspberryPi4Specs} ejecuta Ubuntu~20.04 y \gls{ROS2}~Foxy, orquestando la teleoperación, la percepción (\gls{LIDAR}) y los nodos de SLAM y navegación. Se seleccionó por su compatibilidad oficial con Ubuntu~20.04 y \gls{ROS2}~Foxy, su rendimiento suficiente para SLAM 2D y Nav2 (4 núcleos Cortex-A72 a 1.8~GHz, 8~GB de RAM), su conectividad integrada (USB para el \gls{LIDAR}, Wi-Fi y Ethernet para teleoperación, GPIO) y su amplio ecosistema de documentación \citep{RaspberryPi4Specs, RaspberryPi4Datasheet}.

\subsection{Alternativas consideradas}
\begin{itemize}
  \item \textbf{NVIDIA Jetson Orin Nano}: mayor potencia para IA, pero mayor costo y complejidad.
  \item \textbf{ODROID-N2+}: buen rendimiento, pero ecosistema menos maduro para \gls{ROS2}.
  \item \textbf{BeagleBone Black}: PRUs para tiempo real, pero CPU limitada para SLAM.
\end{itemize}

Dado el alcance del proyecto, la Raspberry Pi~4 es la opción más equilibrada.

%------------------------------------------------------------------
\section{Sensor LiDAR para mapeo y navegación}

El sensor principal de percepción es el RPLidar~A1M8, un escáner láser 2D de 360° por triangulación. La justificación de un \gls{LIDAR}~2D frente a alternativas como cámaras RGB-D, ultrasonido o \gls{LIDAR}~3D para interiores planos de bajo costo se presentó en el Capítulo~\ref{ch:estado_arte}; esta sección detalla sus especificaciones y las consideraciones de integración y montaje relevantes para el prototipo. El sensor se integra de forma nativa en \gls{ROS2} mediante el paquete \texttt{rplidar\_ros}.

\utemFigure{2.Texto/Imag/a1m8.jpg}{0.4\linewidth}{Láser 2D RPLIDAR A1M8.}{NaylampA1M8R5}{fig:lidar}

\paragraph{Especificaciones} según \citep{SlamtecA1Datasheet}
\begin{itemize}
  \item Alcance: 0.15--12~m.
  \item Frecuencia de barrido: 1--10~Hz (típico 5.5~Hz).
  \item Tasa de muestreo: hasta 8000 muestras/s (revisión R5).
  \item Resolución angular: $\le 1^\circ$.
  \item Alimentación: sistema de escaneo a 5~V, consumo total $\sim$450~mA.
  \item Peso: $\sim$170~g.
\end{itemize}

\subsection{Riesgos y mitigación}
\begin{itemize}
  \item \textbf{Superficies transparentes o especulares}: producen lecturas erráticas; se mitigan con capas de inflación conservadoras en el costmap (Capítulo de navegación).
  \item \textbf{Vibraciones}: se amortigua el montaje del sensor y se limitan las aceleraciones bruscas del robot.
\end{itemize}

%------------------------------------------------------------------
\section{Alimentación y distribución de energía}

\subsection{Batería}

Se utiliza una batería Li-ion de 12~V en formato deslizante (compatible con herramientas Makita), elegida frente a alternativas como packs DIY de celdas 18650, LiFePO$_4$ o plomo-ácido por su facilidad de recambio, seguridad y buena relación energía/peso.

\utemFigure{2.Texto/Imag/battery.png}{0.4\linewidth}{Batería Li\mbox{-}ion de 12~V.}{AIMJINBL1021B}{fig:battery}

\subsection{Regulador DC-DC (buck)}

Se usa el módulo DFRobot DFR0205, basado en el chip GS2678, un convertidor reductor conmutado a 350~kHz \parencite{DFRobotPowerModule}. Acepta tensiones de entrada entre 3.6 y 25~V y ofrece una salida fija de 5~V, una ajustable entre 3.3 y 25~V y un paso directo de la entrada, entregando hasta 5~A (25~W) con protecciones por cortocircuito y limitación de corriente. Este regulador deriva el riel de 5~V (Raspberry Pi, \gls{LIDAR}, \gls{Arduino}) desde la batería de 12~V con bajas pérdidas.

\utemFigure{2.Texto/Imag/modulo_dc.jpg}{0.4\linewidth}{Power Module SKU DFR0205.}{DFRobotPowerModule}{fig:modulo_dc}

\subsection{Protecciones e interruptores}

El ramal de potencia sigue la cadena: batería $\rightarrow$ fusible $\rightarrow$ interruptor maestro $\rightarrow$ reguladores/cargas. Se utiliza un fusible rápido tipo \emph{mini blade} de 10~A, que protege el cableado aguas abajo, y un interruptor maestro que secciona el bus de 12~V para cortar la potencia del sistema de forma segura.

\utemFigure{2.Texto/Imag/switch.jpg}{0.4\linewidth}{Interruptor maestro de 16~mm.}{lanboo_16mm_economical_switch}{fig:switch}

\utemFigure{2.Texto/Imag/fuse.jpg}{0.4\linewidth}{Fusible automotriz tipo \emph{blade} Mini APM.}{optifuse_apm}{fig:fuse}

\subsection{Dimensionamiento de cableado y conectores}

La Tabla~\ref{tab:cableado_conectores} resume el dimensionamiento estimado de los ramales de potencia. Los valores de corriente y longitud se ajustan durante el comisionamiento final.

\begin{table}[H]
\centering
\begin{threeparttable}
  \caption{Dimensionamiento de cableado y conectores.}
  \label{tab:cableado_conectores}
  \small
  \begin{tabular}{P{4.2cm} c c c P{3.2cm}}
    \toprule
    \textbf{Ramal de potencia / señal} & \textbf{$I_{max}$ (A)} & \textbf{Long. (m)} & \textbf{$\Delta V$ adm.}\tnote{b} & \textbf{AWG / Conector}\tnote{a} \\
    \midrule
Batería $\rightarrow$ switch principal        & 5.2\tnote{c} & 0.15 & $<0.24$\,V & 14 AWG / XT60    \\
Switch $\rightarrow$ regulador DC-DC (5V)     & 5.0 & 0.15 & $<0.36$\,V & 18 AWG / Bornera \\
Línea 12V $\rightarrow$ driver motores (L298N)& 3.6 & 0.20 & $<0.60$\,V & 18 AWG / Bornera \\
Línea 5V $\rightarrow$ SBC (Raspberry Pi 4)  & 3.0 & 0.25 & $<0.25$\,V & 22 AWG / USB-C   \\
    \bottomrule
  \end{tabular}
  \begin{tablenotes}
    \footnotesize
    \item[a] Los calibres son orientativos; se recomienda cable de silicona multifilar para facilitar el ruteo.
    \item[b] Caída de tensión admisible, calculada como porcentaje del voltaje nominal del ramal.
    \item[c] Pico instantáneo del sistema; protegido por el fusible de 10~A aguas arriba.
  \end{tablenotes}
\end{threeparttable}
\end{table}

%------------------------------------------------------------------
\section{Fabricación y diseño 3D}
\label{sec:fab_3d}

\subsection{Material y equipo}

El chasis y las cubiertas se fabrican mediante impresión 3D FDM en PLA+, por su facilidad de impresión, baja deformación y estabilidad dimensional. La impresión se realiza en una Creality Ender~3 con boquilla de latón de 0.4~mm y filamento de 1.75~mm almacenado con desecante. El PLA+ es adecuado para interiores, aunque se degrada con el calor o la radiación UV.

\subsection{Diseño}

El modelado 3D se realizó en \textit{Fusion~360} de forma paramétrica, organizando el proyecto por componentes (carcasa, tapas de acceso, soporte de \gls{LIDAR}, subchasis y alojamientos electrónicos). Se contemplaron holguras típicas para FDM (0.2--0.4~mm) y se añadieron chaflanes para reducir concentraciones de tensión. Las uniones de uso frecuente emplean insertos roscados por calor (\emph{heat-set}) con tornillería M3/M4, y la orientación de impresión se alineó con los esfuerzos principales para mitigar la anisotropía. Los parámetros de laminado completos se detallan en el Anexo~\ref{anx:impresion3d}.

Las cotas generales se presentan en la Figura~\ref{fig:medidas_robot_2d} y las dimensiones relevantes en la Tabla~\ref{tab:dimensiones_robot}.

\utemFigurePropia{2.Texto/Imag/medidas_robot_2d.pdf}{0.9\linewidth}{Medidas del chasis (vista 2D).}{fig:medidas_robot_2d}

\begin{table}[H]
\centering
\caption{Dimensiones relevantes del robot (según plano de medidas).}
\label{tab:dimensiones_robot}
\small
\begin{threeparttable}
\begin{tabular}{l l p{7cm}}
\toprule
\textbf{Parámetro} & \textbf{Valor} & \textbf{Descripción} \\ 
\midrule
Largo total                  & 246 mm\tnote{a} & Extremo posterior a extremo frontal de la carcasa. \\
Ancho total                  & 230 mm & Lado a lado en el eje de las ruedas. \\
Alto chasis (sin \gls{LIDAR})      & 105 mm & Base a cubierta superior. \\
Alto total (con \gls{LIDAR})       & 153 mm & Base a punto más alto del sensor. \\ \addlinespace
Distancia entre ruedas ($L$) & 188 mm & Centro a centro de ruedas motrices (\textit{track width}). \\
Distancia libre al suelo     & 14 mm  & Base a superficie de apoyo (\textit{ground clearance}). \\
Espesor de pared nominal     & 10 mm  & Paredes de la carcasa impresa en 3D. \\
\bottomrule
\end{tabular}
\begin{tablenotes}
  \footnotesize
  \item[a] Leve discrepancia con el diseño 3D producto de tolerancias de impresión.
\end{tablenotes}
\end{threeparttable}
\end{table}

% ===========================================================================
\section{Galería fotográfica del prototipo}
\label{sec:galeria}
\begin{utemMultiFigurePropia}{Galería fotográfica del prototipo construido.}{fig:galeria_prototipo}

\addutem{0.36}{2.Texto/Imag/fig_prototipo_planta.jpg}{Vista en planta}
\hfill
\addutem{0.42}{2.Texto/Imag/fig_prototipo_lateral_izquierda.jpg}{Vista lateral izquierda}

\vspace{1cm}
\addutem{0.42}{2.Texto/Imag/fig_prototipo_lateral_derecha.jpg}{Vista lateral derecha}
\hfill
\addutem{0.42}{2.Texto/Imag/fig_prototipo_frontal.jpg}{Vista frontal}
\end{utemMultiFigurePropia}

Las vistas isométricas del modelo 3D se incluyen en el Anexo~\ref{anx:renders3d}.

%------------------------------------------------------------------
\section{Parámetros físicos del modelo}
\label{sec:parametros_modelo}

La Tabla~\ref{tab:parametros_modelo} reúne los parámetros físicos del prototipo que alimentan el modelo cinemático y de odometría del Capítulo~\ref{marco_teorico} y la calibración del firmware. Estos valores constituyen la referencia única citada por los capítulos de gemelo digital, firmware y validación.

\begin{table}[H]
\centering
\caption{Parámetros físicos del prototipo empleados en el modelo y el firmware.}
\label{tab:parametros_modelo}
\small
\begin{tabular}{c l l l}
\toprule
\textbf{Símbolo} & \textbf{Parámetro} & \textbf{Valor} & \textbf{Origen} \\
\midrule
$r$   & Radio de rueda                   & 0.0335 m           & Tabla~\ref{tab:ruedas-especificaciones} \\
$C$   & Circunferencia de rueda          & 0.2105 m           & $\pi d$ \\
$L$   & Distancia entre ruedas           & 0.188 m            & Tabla~\ref{tab:dimensiones_robot} \\
$N_e$ & Resolución del encoder (salida)  & 1980 cuentas/vuelta & $11 \times 4 \times 45$ \\
$G$   & Relación de reducción            & $\approx$45        & Motorreductor JGA25-370 \\
$T_s$ & Período del lazo PID             & $1/30$ s           & Firmware del \gls{Arduino} \\
\bottomrule
\end{tabular}
\end{table}

%------------------------------------------------------------------
\section*{Conclusión del capítulo}

En este capítulo se han descrito y justificado los componentes del prototipo: la plataforma diferencial, ruedas, motorreductores, encoders, driver L298N, microcontrolador \gls{Arduino}~Nano, computador Raspberry Pi~4, sensor \gls{LIDAR} RPLidar~A1 y el sistema de alimentación, así como la fabricación del chasis mediante impresión 3D. Cada selección se fundamentó en criterios de costo, disponibilidad, compatibilidad con \gls{ROS2} y adecuación a los requisitos del proyecto. Los parámetros físicos resultantes, consolidados en la Tabla~\ref{tab:parametros_modelo}, alimentan los capítulos siguientes, donde se construyen el gemelo digital, el firmware de bajo nivel y la arquitectura de control y navegación.